% Ad Atticum 10.1 (ad-atticum-10-01) --- English translation
% Translated by Claude (Anthropic) from the Latin (Perseus).
% Style guide: STYLE.md.

\ciceroLetterOpener{CICERO ATTICO salutem}

\ciceroSection{1} On the third day before the Nones, when I had reached my brother's place at Laterium, I got your letter and breathed a little easier --- which has not happened to me since this ruin. For I count it of the highest importance that you should approve our firmness of mind and what we have done. As for your writing that our Sextus approves it too, I take such delight in this that I think myself confirmed as it were by the judgement of his father, the one man to whom I always gave the greatest weight. I often remember what he said to me long ago, on that famous Nones of December, when I asked: ``Well, Sextus, what now?'' His answer was: ``Not without a struggle and inglorious, but having done some great thing for men of after-time to learn of'' \textit{[Greek: m\=e m\'an asp\=oud\'i ge ka\`i akle\=eos, all\`a m\'ega rh\'exas ti ka\`i essomenoisi puth\'esthai]}. His authority, then, lives on for me, and his son, the very image of him, carries the same weight with me as he did. Please give him my warmest greetings.

\ciceroSection{2} You put off your own counsel not to a remote time (for by now I take it that bought-up peacemaker has finished his speech, by now some business has been transacted in a gathering of senators --- senate I do not call it), and yet you keep my mind in suspense; but less so, because I do not doubt what you think we ought to do. For you, who write that Flavius is being given a legion and Sicily, and that this is already happening --- what crimes do you reckon are some of them being prepared and contemplated already, and some of them yet to be improvised? I, then, will set aside the law of Solon, your countryman and, I take it, mine too, who made it a capital offence for a man in time of civil strife to belong to neither side; and, unless you judge otherwise, I will be absent both from this side and from that. But of the two the second is the more settled with me; I will not anticipate, however. I shall wait for your counsel and the letter --- unless you have already sent another, the one I asked you to give to Cephalio.

\ciceroSection{3} As for your writing that, not on any external report but by your own judgement, you think I shall be dragged in if there is any question of peace --- it does not so much as come into my mind what action there could possibly be about peace, when for him it is fixed beyond doubt, if he can, to strip Pompey of army and province; unless perhaps that money-man can prevail on him to keep quiet while ambassadors go and return. I see nothing to hope for, nothing that I can now even think possible. And yet this very thing belongs to an honest man: †let it be a great matter† among the most properly political questions \textit{[Greek: t\^on politik\=ot\'at\=on skemm\'at\=on]} --- whether one should attend the council of a tyrant when he is going to deliberate on some good business. So if anything of the kind happens and I am summoned (which I do not for my part believe; for what should I have to say about peace? I have said my say, and he himself rejected it firmly) --- but still, if anything does happen, by all means write me what you think I ought to do. For nothing has yet befallen me on which the deliberation was greater. As for Trebatius, an honest man and citizen, I am glad you took pleasure in his words, and that frequent exclamation of yours, ``Splendid!'' \textit{[Greek: ekph\=on\=esis hup\'ereu]}, is the only thing that has given me pleasure up to now. I am keenly waiting for your letter; which by now, I take it, has been dispatched.

\ciceroSection{4} You held with Sextus the same gravity that you prescribe to me. Your Celer is more eloquent than wise. What you heard from Tullia about the young men is true. The business about Mucianus that you write of seems to me not so harsh in fact as in word. This is the wandering \textit{[Greek: \'al\=e]} we are now in --- it is as good as death. For either I had to play my part as citizen \textit{[Greek: politeut\'eon]} freely among the wicked, or, with some danger, with the honest. Either we follow the recklessness of the honest, or harry the audacity of the unprincipled. Either is dangerous; but this thing we are doing is disgraceful and not even safe. As for that man who sent his son to Brundisium about peace --- about peace I think the same as you, that the pretence is open and that war is being prepared with the keenest effort --- I do not believe I shall be sent as envoy by him; and so far, as I hoped, no mention has been made of me. The less, then, I need to write or even to think about what I shall do, if it does come about that I am chosen.
